% Preamble for the reMarkable-sized PDF exports built by the Makefile.
%   reMarkable 1 — 1404x1872 px @ 226 dpi -> 447.29 x 596.39 bp (≈ 157 x 210 mm)
% The page box matches the panel aspect ratio so the reader shows the page
% full-screen with no letterboxing and no pinch-zooming.
\usepackage[paperwidth=447.29bp,paperheight=596.39bp,%
            top=11mm,bottom=13mm,left=11mm,right=11mm]{geometry}
\usepackage{amsmath}
\usepackage{amssymb}
\usepackage{etoolbox}
\usepackage{fvextra}
\usepackage{newunicodechar}
\usepackage{pifont}

% FreeSerif/FreeMono lack these glyphs in some faces (bold, italic); route them
% through math mode so every weight renders identically.
\newunicodechar{ₙ}{\textsubscript{n}}
\newunicodechar{ℕ}{\ensuremath{\mathbb{N}}}
\newunicodechar{ℚ}{\ensuremath{\mathbb{Q}}}
\newunicodechar{ℝ}{\ensuremath{\mathbb{R}}}
\newunicodechar{ℤ}{\ensuremath{\mathbb{Z}}}
\newunicodechar{★}{\ding{72}}
\newunicodechar{⟺}{\ensuremath{\Longleftrightarrow}}
\newunicodechar{✓}{\ding{51}}
\newunicodechar{✗}{\ding{55}}

% unicode-math maps \setminus to U+29F5, which Latin Modern Math lacks.
\AtBeginDocument{\renewcommand{\setminus}{\mathbin{\backslash}}}

% Wrap long code lines instead of running off the narrow page.
\fvset{breaklines=true,breakanywhere=true,fontsize=\small}

% Tables are set a step smaller than the body text, with a little extra row
% height so that fractions in cells do not touch the neighbouring rows.
\AtBeginEnvironment{longtable}{\footnotesize\renewcommand{\arraystretch}{1.4}}

% Long derivations may break across pages; ragged bottom avoids stretched gaps.
\allowdisplaybreaks
\raggedbottom
\setlength{\emergencystretch}{3em}

% Shrink over-wide display math to the column instead of letting it run
% into the margin. On these narrow pages a handful of one-line derivations
% (Taylor polynomials, telescoping sums) are wider than the text block.
% All display math in these sources is single-line, so boxing it is safe.
\usepackage{graphicx}
\newsavebox{\fitmathbox}
\AtBeginDocument{%
  \renewcommand{\[}{%
    \par\addvspace{\abovedisplayskip}%
    \begingroup\begin{lrbox}{\fitmathbox}$\displaystyle}%
  \renewcommand{\]}{$\end{lrbox}%
    \noindent\makebox[\linewidth]{%
      \ifdim\wd\fitmathbox>\linewidth
        \resizebox{\linewidth}{!}{\usebox\fitmathbox}%
      \else
        \usebox\fitmathbox
      \fi}%
    \par\addvspace{\belowdisplayskip}\endgroup}%
}

% Let inline math break after commas. Long enumerations such as
% $\mathbb{Z} = \{\dots, -2, -1, 0, 1, 2, \dots\}$ are otherwise a single
% unbreakable chunk and stick out past the edge of a narrow screen.
\Umathcharnumdef\origcomma=\Umathcodenum`\,
\begingroup\catcode`\,=\active
\gdef,{\origcomma\allowbreak}
\endgroup
\AtBeginDocument{\mathcode`\,="8000 }

% Definitions and theorems are set as blockquotes; the default 2.5em indent
% on both sides is a lot of width to give up on a small screen.
\AtBeginDocument{%
  \renewenvironment{quote}%
    {\list{}{\leftmargin=1em\rightmargin=0pt}\item\relax}%
    {\endlist}%
}

% Light grey, page-breakable box for proofs that were not given in the
% lecture. Markdown `::: kiegeszites` divs are turned into this environment by
% kiegeszites.lua.
\usepackage{xcolor}
\usepackage{framed}
\definecolor{kiegeszitesszin}{gray}{0.91}
\newenvironment{kiegeszites}%
  {\def\FrameCommand{\fboxsep=6pt\colorbox{kiegeszitesszin}}%
   \MakeFramed{\advance\hsize-\width\FrameRestore}}%
  {\endMakeFramed}
